
\section{Robustness and Extensions}
\label{sec:robustness}


We subject the baseline to five groups of checks:
(i)~FL-definition variants (IQR multiplier, win-rate cutoff, temporal
window); (ii)~sample and specification variants (unrestricted sample,
homogeneous subsamples, item$\times$year FE, matching);
(iii)~alternative clustering; (iv)~sensitivity to unobservables
(Cinelli--Hazlett); and (v)~staggered DiD.
Table~\ref{tab:robustness_summary} collects the point estimates, with
an extra column spelling out what each estimator assumes.


\begin{table}[htbp]
\centering
\caption{Robustness summary: FL price coefficient across checks}
\label{tab:robustness_summary}
\small
\begin{tabular}{lcccll}
\hline\hline
Check & Coeff.\ & SE & $N$ & Assumption & Note \\
\hline
\multicolumn{6}{l}{\textit{Baseline and matching}} \\
OLS (item + year + PBU FE) & 0.064 & 0.020 & 1{,}654{,}401 & Sel.\ on obs.\ + FE & Baseline \\
CEM & 0.077 & 0.024 & 969{,}751 & Sel.\ on obs. & Sec.~\ref{sec:results_ols} \\
IPW & 0.055 & 0.021 & 830{,}194 & Sel.\ on obs. & Sec.~\ref{sec:results_ols} \\
Cross-fit average & 0.036 & 0.019 & 1{,}654{,}401 & Sel.\ on obs.\ + FE; no mech.\ link & Out-of-sample \\
\hline
\multicolumn{6}{l}{\textit{FL definition variants}} \\
IQR 1.0$\times$ & 0.079 & 0.023 & 1{,}654{,}401 & Sel.\ on obs.\ + FE & 3{,}442 FL firms \\
IQR 1.5$\times$ & \multicolumn{5}{l}{\textit{(= OLS baseline above)}} \\
IQR 2.0$\times$ & 0.060 & 0.022 & 1{,}654{,}401 & Sel.\ on obs.\ + FE & 2{,}093 FL firms \\
IQR 3.0$\times$ & 0.050 & 0.024 & 1{,}654{,}401 & Sel.\ on obs.\ + FE & 1{,}456 FL firms \\
\hline
\multicolumn{6}{l}{\textit{Specification variants}} \\
Item$\times$year FE & 0.074 & 0.022 & 1{,}654{,}401 & Sel.\ on obs.\ + FE & Tighter controls \\
Two-way clustering & 0.064 & 0.024 & 1{,}654{,}401 & Sel.\ on obs.\ + FE & Item + PBU \\
\hline
\multicolumn{6}{l}{\textit{Alternative treatments and diagnostics}} \\
Continuous (log max AL) & 0.022 & 0.005 & 1{,}654{,}401 & Sel.\ on obs.\ + FE & Per log-point \\
Decomposition (odd-yr FL, full sample) & 0.043 & 0.019 & 1{,}654{,}401 & Sel.\ on obs.\ + FE & Intermediate \\
Horse-race (FL + Imhof CV) & 0.084 & 0.023 & 1{,}654{,}401 & Sel.\ on obs.\ + FE & Suppression \\
IV (leave-one-out) & 0.194 & 0.077 & 1{,}654{,}401 & Excl.\ restriction & Measurement error \\
\hline
\multicolumn{6}{l}{\textit{Staggered DiD and coefficient stability}} \\
Stacked DiD (Cengiz et al.) & $-$0.006 & 0.014 & 715{,}116 & Parallel trends & Pre-trends violated \\
Oster $\hat{\delta}$ & \multicolumn{2}{c}{degenerate} & 1{,}654{,}401 & Proportional selection & $R^2$ gap $\approx 0$ \\
\hline\hline
\end{tabular}
\begin{tablenotes}
\small
\item \textit{Notes:} ``Sel.\ on obs.\ + FE'' = selection on observables absorbed by item, year, and PBU fixed effects. Cross-fit further removes any mechanical link between the FL rule and the price outcome. The IV is a measurement-error diagnostic; exclusion-restriction concerns (Section~\ref{sec:results_ols}) keep it off the headline range.
\end{tablenotes}
\end{table}


\paragraph{Threshold sensitivity.}
The coefficient decreases monotonically with stricter thresholds
(0.079 at 1.0$\times$, 0.064 baseline, 0.060 at 2.0$\times$, 0.050
at 3.0$\times$), most likely reflecting classification noise as the
treated group shrinks. A two-dimensional sensitivity analysis varying
both multiplier and win-rate cutoff produces significant coefficients
across all 36 cells
(Figure~\ref{fig:threshold_heatmap}, Appendix~\ref{app:robustness}).

\paragraph{Sensitivity to unobservables.}
\citet{cinelli2020making} sensitivity analysis yields $RV_{q=1} =
17.5\%$: an unobserved confounder would need to explain 17.5\% of
residual variation in \emph{both} FL presence and log prices to nullify
the coefficient---stringent given the item, year, and PBU fixed-effects
structure. Full sensitivity plots are in
Appendix~\ref{app:robustness}.


\paragraph{Staggered difference-in-differences.}
\citet{callaway2021difference} yields ATT~$= 0.014$ (SE~$= 0.039$);
stacked DiD \citep{cengiz2019effect} gives $-0.006$ (SE~$= 0.014$;
Table~\ref{tab:stacked_did}, Appendix~\ref{app:did}). Positive
pre-trend coefficients at $t = -2$ and $t = -3$ preclude a causal
reading: PBUs receiving FL entry \emph{already had higher prices},
consistent with strategic market selection but equally compatible
with within-market confounding. The pre-trends reinforce the paper's
non-causal framing and underscore that within-item conditioning does
not eliminate time-varying within-market confounders.
\citet{rambachan2023more} sensitivity analysis is in
Appendix~\ref{app:did}.

\paragraph{Coefficient stability (Oster).}
The \citet{oster2019unobservable} bound is degenerate
($\hat{\delta} = 261.6$) because PBU FE barely move $R^2$
(0.879 to 0.886)---a design strength, not fragility. The
\citet{cinelli2020making} framework ($RV = 17.5\%$) is the
appropriate sensitivity metric
(Table~\ref{tab:oster_delta}, Appendix~\ref{app:robustness}).


\paragraph{Illustrative welfare.}
Appendix~\ref{app:extensions} reports back-of-the-envelope welfare
calculations applying the price coefficients to FL-present spending,
under an assumption of full causality.


\paragraph{Oversight heterogeneity.}
The conditional FL--price association varies with PBU size: 0.214 in
the smallest quartile, 0.098 in Q2, 0.045 in Q3, and 0.017 in Q4
(Table~\ref{tab:regime_oversight}, Appendix~\ref{app:robustness}).
The 12.5$\times$ gradient is consistent with the framework's
comparative static ($\partial m^*/\partial \theta_k < 0$): where
oversight capacity is lower, the FL--price association is larger.
The pattern is suggestive from a policy standpoint---even modest
investments in procurement monitoring may compress the scope for
anomalous entry---though not uniquely diagnostic, since smaller PBUs
also have thinner fixed-effect cells that could mechanically inflate
the coefficient.%
\footnote{The monotonic gradient across all four quartiles weighs
against a pure noise story, but we cannot fully rule out that thinner
item--PBU cells in smaller PBUs drive part of the pattern.}
Threshold stability across IQR multipliers and a comparison between
the IQR rule and model-based classifiers are in
Appendix~\ref{app:robustness}.
